\documentclass[main.tex]{subfiles}
\begin{document}
\section{Attivazione del Payload}\label{sec:cap 4}

Dopo aver individuato correttamente l’offset, calcolato l’indirizzo del buffer e impostato il return address, possiamo procedere con l’ultima fase dell’attacco: l’esecuzione dell’exploit.
Questa fase consiste nell’invio al programma vulnerabile di un payload costruito ad hoc, il quale combina la shellcode, un’area di padding per raggiungere il return address e infine l’indirizzo di ritorno scelto. L’obiettivo è dirottare il flusso di esecuzione verso la shellcode, ottenendo così l’esecuzione arbitraria di comandi.\\
Il payload viene scritto su un file binario e poi fornito come input al programma, facendo crashare il flusso normale per aprire una shell con i privilegi dell’utente che esegue il binario.\\
\noindent Eseguendo da terminale il comando \texttt{(cat payload.bin; cat) | ./vuln  } si aprirà la nostra shell ed avremo un output di questo tipo:

\begin{figure}[h]
    \centering
    \includegraphics[width=0.5\textwidth]{figures/payload.png}
    \caption{Output della compilazione del file exploit.py da terminale}
    \label{fig:esempio}
\end{figure}

\noindent Ora, possiamo eseguire comandi quali id, whoami, ls -l, ed interagire con la nostra shell.

\begin{figure}[h]
    \centering
    \includegraphics[width=0.8\textwidth]{figures/shell.png}
    \caption{Esecuzione della shell}
    \label{fig:esempio}
\end{figure}

\noindent Il nostro attacco è riuscito, e siamo stati in grado anche di ottenere i permessi root con la nostra shell.





\end{document}